\chapter{Hashtables}\label{chap:hashtables}

Hashtables are a common implementation of the map ADT. They use a hash function to map keys to indices in an array (buckets).

\section{Design rationale}

Our hashtable implementation focuses on performance and flexibility. We support both open addressing and closed addressing strategies (detailed in subsequent chapters).

Key design decisions include:
\begin{itemize}
    \item \textbf{Resizing:} To maintain $O(1)$ average time complexity, the hashtable must resize when the load factor (ratio of elements to buckets) exceeds a certain threshold. We implement automatic resizing (rehashing) to a larger size (typically doubling the capacity).
    \item \textbf{Hashing:} The hash function is a parameter provided by the user. This allows the user to choose a hash function appropriate for their specific key type and distribution.
    \item \textbf{Genericity:} Like other data structures in this library, hashtables are implemented using macros to achieve type safety and performance (avoiding \texttt{void*} overhead where possible) while remaining generic.
\end{itemize}

\section{Hashtable interface}

The hashtable interface conforms to the general map interface described in Chapter~\ref{chap:maps}. It provides specific implementations for \texttt{put}, \texttt{get}, \texttt{remove}, etc., based on the chosen collision resolution strategy.
